% Law of total variance, drawn as three conditional densities of Y at
% three values of the conditioning variable X. The within-group spread
% is the average conditional variance; the between-group spread is the
% variance of the three conditional means. Total variance is their sum.
% Densities are stylised Gaussians at fixed centres and widths.
\begin{figure}[htbp]
\centering
\begin{tikzpicture}[x=1cm, y=1cm, font=\small]

% Horizontal axis (Y)
\draw[->, thick] (0, 0) -- (11.4, 0) node[right] {$Y$};
\draw[->, thick] (0, 0) -- (0, 3.6) node[above] {conditional density};

% Three conditional densities of Y | X = x_i, centred at 2.4, 5.4, 8.4
% with the same within-group sigma ~ 0.7 (equal widths).
% Height factor scales 1/(sigma sqrt(2pi)) into plot units.
% Group 1 (X = x_1), blue
\draw[thick, engblue, smooth]
  plot[domain=0.4:4.4, samples=70]
  (\x, {2.3*exp(-0.5*pow((\x-2.4)/0.7,2))});
\draw[dashed, engblue!70] (2.4, 0) -- (2.4, 2.3);
\node[anchor=north, engblue, font=\scriptsize] at (2.4, -0.05)
  {$\E[Y\mid X{=}x_1]$};

% Group 2 (X = x_2), amber
\draw[thick, engamber, smooth]
  plot[domain=3.4:7.4, samples=70]
  (\x, {2.3*exp(-0.5*pow((\x-5.4)/0.7,2))});
\draw[dashed, engamber!80] (5.4, 0) -- (5.4, 2.3);
\node[anchor=north, engamber!80!black, font=\scriptsize] at (5.4, -0.05)
  {$\E[Y\mid X{=}x_2]$};

% Group 3 (X = x_3), red
\draw[thick, engred, smooth]
  plot[domain=6.4:10.4, samples=70]
  (\x, {2.3*exp(-0.5*pow((\x-8.4)/0.7,2))});
\draw[dashed, engred!70] (8.4, 0) -- (8.4, 2.3);
\node[anchor=north, engred, font=\scriptsize] at (8.4, -0.05)
  {$\E[Y\mid X{=}x_3]$};

% Within-group arrow (one sigma on group 2)
\draw[<->, thick, enggray] (4.7, 1.35) -- (6.1, 1.35);
\node[anchor=south, enggray, font=\scriptsize] at (5.4, 1.35)
  {within: $\E[\Var(Y\mid X)]$};

% Between-group arrow (spread of the three means)
\draw[<->, thick, engblack] (2.4, 2.95) -- (8.4, 2.95);
\node[anchor=south, engblack, font=\scriptsize] at (5.4, 2.95)
  {between: $\Var(\E[Y\mid X])$};

\end{tikzpicture}
\caption[The law of total variance]{The variance decomposition
$\Var(Y)=\E[\Var(Y\mid X)]+\Var(\E[Y\mid X])$ read off three
conditional densities of $Y$, one for each value of the conditioning
variable $X$. The \emph{within-group} term (grey) is the average spread
of $Y$ inside a group, the residual variance left after conditioning.
The \emph{between-group} term (black) is the spread of the three group
means, the variance explained by $X$. Their sum is the total variance
of $Y$. The decomposition is the probabilistic skeleton of the analysis
of variance, of variance-component models in metrology, and of the
bias-variance accounting of the machine-learning chapters. Source:
schematic by the authors; the compound-sum worked example of section
13.3 applies the same identity to a random part count.}
\label{fig:vol02ch13:total-variance}
\end{figure}
